% !TeX root = ../main.tex
% Rebuilt from: ch08-numerical-modelling.tex, advanced-numerical.tex
\chapter{Numerical Modelling}
\label{ch:numerical-modelling}

Numerical analysis removes the geometric and constitutive restrictions that
limit closed-form solutions, at the cost of requiring parameters and
assumptions that the analyst must supply and defend. This chapter treats the
two applications that dominate underground practice: prediction of
tunnelling-induced ground movement, and analysis of excavation and support
response in rock. It then considers the methods that lie beyond continuum
analysis and the coupled formulations that some problems require, sets out
the calibration and quality assurance on which any reliance upon numerical
output depends, and closes with the surrogate models that make repeated
evaluation affordable.

\section{Ground movement and model set-up}
\label{sec:ground-movement-and-model-set-up}

\subsection{The Gaussian settlement trough}
\label{sec:gaussian-trough}

Surface settlement above a tunnel is predicted twice in most projects, first
empirically, before any model exists, and then numerically. The numerical
estimate does not supersede the empirical one, which remains the check
against which the model is judged. A model that cannot reproduce it needs
an explanation before its other outputs are relied upon.

The transverse settlement trough above a single tunnel is well described by a
Gaussian curve~\citep{peck1969deep} (Figure~\ref{fig:gaussian-trough}),
\[
S_y = S_{\max} \, \exp\!\left(-\frac{y^2}{2i^2}\right),
\]
in which $S_y$ is the settlement at transverse distance $y$ from the
centreline, $S_{\max}$ the settlement over the axis, and $i$ the distance
from the centreline to the point of inflection. The trough width follows the
depth to the tunnel axis, $z_0$, through
\[
i = K \, z_0 ,
\]
with $K$ depending on the ground rather than on the tunnel. Table~\ref{tab:k_values}
gives the customary ranges.

\begin{table}[htbp]
\caption{Trough width parameter $K$, after \citet{peck1969deep} and \citet{mair1996settlement}.}
\label{tab:k_values}
\centering
\small
\begin{tabular}{lc}
\toprule
\textbf{Ground type} & \textbf{$K$} \\
\midrule
Soft clay & 0.5--0.7 \\
Stiff clay & 0.4--0.5 \\
Granular soil (sand and gravel) & 0.25--0.45 \\
Weak rock & 0.25--0.4 \\
\bottomrule
\end{tabular}
\end{table}

\begin{figure}[htbp]
\centering
\datafigure[\textwidth]{ch17/gaussian-trough}
\caption{Settlement trough and horizontal strain above a 6~m tunnel at 20~m depth, $K = 0.5$ and $V_L = 1$ per cent. Dotted lines mark the points of inflection at $\pm i$; the strain is compressive between them and tensile outside.}
\label{fig:gaussian-trough}
\end{figure}

The magnitude follows from volume loss (Table~\ref{tab:vl_values}). The volume of the trough per metre of
tunnel, $V_s$, is expressed as a percentage $V_L$ of the excavated area, with
$D$ the excavated diameter, which for a shielded machine is the bored
diameter and not the internal diameter of the finished lining,
\[
V_s = \frac{V_L}{100} \times \frac{\pi D^2}{4},
\]
and equating this to the integral of the Gaussian gives the maximum
settlement,
\[
S_{\max} = \frac{V_s}{i \sqrt{2\pi}} .
\]

\noindent\textbf{A note on units.}\enspace The factor of 100 appears only in the
expression for $V_s$. A volume loss quoted as 0.5\,per cent
enters that expression as 0.005; $S_{\max}$ then emerges in metres and is
converted to millimetres at the end. Applying the percentage twice, or not at
all, is the commonest arithmetic error in this calculation and puts the answer
out by two orders of magnitude, which makes it easy to detect.

\begin{table}[htbp]
\caption{Typical volume loss by excavation method.}
\label{tab:vl_values}
\centering
\small
\begin{tabular}{lc}
\toprule
\textbf{Excavation method} & \textbf{Typical $V_L$ (\%)} \\
\midrule
Closed-face TBM in clay or sand & 0.5--1.5 \\
Open-face TBM in rock & 0.5--2.0 \\
Sprayed concrete lining in soft ground & 1.0--3.0 \\
\bottomrule
\end{tabular}
\end{table}

Volume loss has several sources, and treating it as a single quantity obscures
where control is available. In shielded tunnelling it accumulates from five
sources~\citep{mair1996settlement} (Figure~\ref{fig:volume-loss-sources}): stress relief at the face, over-cutting as
the shield passes, closure of the tail void before grout takes effect,
deflection of the lining under ground load, and consolidation driven by the
change in pore pressure. The first three are controlled by the machine and by
the operator, the fourth by the lining design, and the fifth by the
hydrogeology. A volume loss target written into a specification therefore
covers all five, and is meaningful only if the contract makes clear
which party controls each.

\begin{figure}[htbp]
\centering
\datafigure[\textwidth]{ch17/volume-loss-sources}
\caption{Sources of volume loss in shielded tunnelling (schematic, gaps exaggerated): 1 face, 2 over-cutting along the shield, 3 tail void closing before the grout acts, 4 lining deflection (cross-section, right), 5 consolidation. The shaded band at the crown is ground lost into the cut profile; the dashed line is the settled surface.}
\label{fig:volume-loss-sources}
\end{figure}

\subsection{Model inputs and ground idealisation}
\label{sec:model-inputs-and-geological-idealisation}

A numerical model of a shallow tunnel is built from a deliberately simplified
ground profile. The simplification is justified, because a
model carrying every unit identified in the investigation is slow, hard to
interpret and no more reliable than the parameters it is given. The units
retained should be those that govern the answer, and the
simplification should be recorded.

For settlement analysis the governing distinction is usually between a
compressible surface unit and the stiffer material in which the tunnel sits.
Table~\ref{tab:props} gives a two-unit idealisation of the kind used for a
shallow urban alignment, with a granular soil over a weak sedimentary rock.

\begin{table}[htbp]
\caption{A two-unit ground idealisation for shallow tunnel analysis, with a
Mohr--Coulomb constitutive model. Values of this order are typical of a
weathered urban profile and are given to fix the worked examples that follow;
project values come from the ground model.}
\label{tab:props}
\centering
\small
\begin{tabular}{l l c c c c c c c}
\toprule
\textbf{Unit} & \textbf{Description} & $\gamma$ & $c'$ & $\phi'$ & $\psi$ & $E_i$ & $\nu$ & $K_0$ \\
 & & (kN/m\textsuperscript{3}) & (kPa) & (\textdegree) & (\textdegree) & (MPa) & & \\
\midrule
Upper & Granular soil, surface layer & 20 & 7 & 30 & 2 & 50 & 0.3 & 0.7 \\
Lower & Weak sedimentary rock, host & 21 & 20 & 28 & 2 & 150 & 0.3 & 0.75 \\
\bottomrule
\end{tabular}
\end{table}

Two aspects of the input deserve more care than they usually receive. The first
is the stress ratio. Where the horizontal stress differs between the in-plane
and out-of-plane directions (a common condition in bedded or tectonised
rock), the plane-strain model must be given both, since the out-of-plane
stress affects yielding through the intermediate principal stress even though
no displacement occurs in that direction. Adopting a single isotropic $K_0$
where the field is anisotropic is a defensible simplification, but it should be
stated, and its effect on the computed settlement is worth checking with a
separate run.

The second is provenance. Parameters published in an approved environmental
statement or a geotechnical baseline report carry contractual weight, and
departing from them requires a reason recorded in the design file. Where the
profile is simplified, by merging units or by adopting a single $K_0$, the
simplification and its justification belong in the same place.

The support is idealised in the same spirit. Table~\ref{tab:lining} gives
properties for a continuous concrete ring used to represent a bolted segmental
lining in two dimensions. The substitution is legitimate because a
plane-strain analysis represents only axial and bending stiffness and a
strength limit. It does not represent the joints, the bolts or the ring-build
sequence, so the label attached to the material matters less than the
stiffness and capacity entered. The articulation at the joints, and the
reduction in effective bending stiffness that follows from it, are omitted.
They are the subject of Chapter~\ref{ch:linings} and are the reason a
two-dimensional lining analysis tends to overestimate moment.

\begin{table}[htbp]
\caption{Lining properties for a two-dimensional analysis, representing a
segmental ring by an equivalent continuous section.}
\label{tab:lining}
\centering
\small
\begin{tabular}{l l l}
\toprule
\textbf{Element} & \textbf{Type} & \textbf{Properties} \\
\midrule
Concrete ring & Beam, elastic--plastic & $E = 40$~GPa, $\nu = 0.25$, $\gamma = 25$~kN/m\textsuperscript{3} \\
 & & $\sigma_{c,\text{peak}} = 50$~MPa, $\sigma_{t,\text{peak}} = 4$~MPa \\
 & & $\sigma_{c,\text{res}} = 10$~MPa, $\sigma_{t,\text{res}} = 0$~MPa \\
\bottomrule
\end{tabular}
\end{table}

Where the tunnel straddles a contact between two units, the position of that
contact becomes a design assumption. A crown sitting in the
weaker unit and an invert in the stronger produce a different failure
mechanism, and a different trough, from a tunnel wholly within either. The
assumed contact depth should be stated, its basis given, and its influence
tested by moving it within the range the boreholes allow. Here
a sensitivity run costs an hour and resolves a question that would otherwise
remain open throughout the design.

\subsection{Worked examples}
\label{sec:worked-examples}

\paragraph{Initial stress.}
The profile is taken as dry, so that total and effective stresses coincide
and the stress ratio may be applied directly; with a water table above axis
level the ratio applies to the effective stress and the pore pressure is
added afterwards. For a tunnel axis at $z_0 = 25$~m with the contact between the two units of
Table~\ref{tab:props} taken at 2.5~m depth, the total vertical stress at axis
level is
\[
\sigma_v = (20)(2.5) + (21)(25 - 2.5) = 50 + 472.5 = 522.5\ \text{kPa} .
\]
With an in-plane stress ratio of 0.75 and an out-of-plane ratio of 0.9, the
horizontal stresses are
\[
\sigma_h^{\text{in-plane}} = 0.75 \times 522.5 = 391.9\ \text{kPa},
\qquad
\sigma_h^{\text{out-of-plane}} = 0.9 \times 522.5 = 470.3\ \text{kPa} .
\]
A model given a gravity field and the correct unit weights computes these
itself. Working them by hand once is worth the few minutes, because it is the
only practical check that the field stress has been configured as intended.

\paragraph{Settlement trough.}
A separate illustration, on a different geometry from the alignment above,
shows the influence of depth and of ground type. Consider a tunnel of
diameter $D = 6$~m at axis depths of 5~m and 20~m, in
ground taken first as stiff clay ($K = 0.5$) and then as rock ($K = 0.25$),
with $V_L = 1$\,per cent throughout. The trough area is
\[
V_s = 0.01 \times \frac{\pi \times 6^2}{4} = 0.283\ \text{m}^2 ,
\]
and for the shallow tunnel in stiff clay, $i = 0.5 \times 5 = 2.5$~m, so
\[
S_{\max} = \frac{0.283}{2.5 \sqrt{2\pi}} = 0.0451\ \text{m} = 45.1\ \text{mm} .
\]
For the deep tunnel in rock, $i = 0.25 \times 20 = 5.0$~m and
$S_{\max} = 22.6$~mm. Table~\ref{tab:teaching_example} collects the four
combinations, and Figure~\ref{fig:trough-cases} draws them.

\begin{table}[htbp]
\caption{Settlement trough for $D = 6$~m and $V_L = 1$\,per cent. Doubling the
axis depth halves $S_{\max}$ because $i$ doubles; assuming rock rather
than clay doubles it because $i$ halves.}
\label{tab:teaching_example}
\centering
\small
\begin{tabular}{l l c c c}
\toprule
\textbf{Axis depth} & \textbf{Ground assumption} & $i$ (m) & $S_{\max}$ (mm) & \textbf{Trough} \\
\midrule
5~m & Stiff clay ($K=0.5$) & 2.5 & 45.1 & Wider, shallower \\
5~m & Rock ($K=0.25$) & 1.25 & 90.2 & Narrower, deeper \\
20~m & Stiff clay ($K=0.5$) & 10.0 & 11.3 & Wide, very shallow \\
20~m & Rock ($K=0.25$) & 5.0 & 22.6 & Narrower, shallower \\
\bottomrule
\end{tabular}
\end{table}

\begin{figure}[htbp]
\centering
\datafigure[\textwidth]{ch17/trough-cases}
\caption{The four troughs of Table~\ref{tab:teaching_example}, for $D = 6$~m and $V_L = 1$ per cent. All four enclose the same area; the trough width parameter alone sets their shape. At $z_0 = 5$~m the cover is below $1.5D$ and the Gaussian is indicative only.}
\label{fig:trough-cases}
\end{figure}

Two conclusions follow, and they justify making the calculation
before any model is built. First, the trough width parameter, not the volume loss,
dominates the shape: halving $K$ at constant $V_L$ doubles the peak settlement
and narrows the trough, which moves the damage assessment of
Chapter~\ref{ch:settlement-damage} from a wide zone of small differential
movement to a narrow zone of large. Second, the ground assumption has as much
influence as the tunnelling assumption, so where the ground is
intermediate (weathered rock, a soil--rock interface near crown level),
both bounds should be carried forward rather than a single central value.
The combinations multiply quickly once both ground bounds and a range of
$V_L$ are carried, and the trough for any set of $D$, $z_0$, $K$ and $V_L$ is
computed in \companionname\ (Appendix~\ref{app:software}).

\paragraph{Where the Gaussian fails.}
The Gaussian form is an empirical description drawn from case histories of
tunnels with cover of at least one diameter. Where $z_0 < 1.5D$, so that the cover
above the crown is less than one diameter, the mechanism approaches that of a
shallow trapdoor rather than a deep cavity: the trough becomes narrower and
more angular than the curve predicts, and the peak is underestimated. Results
in that range are indicative only, and the limitation should be stated with
the results.

\section{Finite element practice}
\label{sec:finite-element-practice}

\subsection{Building and interpreting a continuum model}
\label{sec:finite-element-modelling}

The finite element method divides the ground into elements connected at nodes,
assembles the element stiffnesses into a global system, and solves for nodal
displacement under the applied loading. For two-dimensional tunnel analysis the
elements are usually six-noded triangles under plane strain. The discussion here
does not depend on which program performs that assembly. The decisions that
determine the value of the answer are made before the solver is
run and after it returns.

The sequence is common to every package. The domain is drawn, with the ground
surface, the significant unit boundaries and the excavation profile. Materials
are assigned, with a constitutive model chosen for each. The mesh is generated
and graded. Boundary conditions are applied
(Figure~\ref{fig:model-domain}): vertical sides restrained
horizontally and free vertically, the base restrained in both directions, the
ground surface free. An initial stress field is established under gravity with
the appropriate stress ratios, and equilibrium checked before anything is
excavated. The excavation is then performed in stages, with support activated
at the stage that represents its actual installation point. Finally the results
are extracted: surface displacement along the ground surface, displacement
contours, structural forces in the lining, and the extent of yielded material.

\begin{figure}[htbp]
\centering
\datafigure[0.78\textwidth]{ch17/model-domain}
\caption{Domain and boundary conditions of a plane-strain tunnel model (schematic): sides on rollers at least $5D$ from the centreline, base fixed at least $5D$ below the invert, surface free, upper unit shaded. The mesh is graded outwards from the excavation, finest within $2D$ of it; it is drawn coarser than the sizes recommended in the text.}
\label{fig:model-domain}
\end{figure}

Two of those steps repay attention. The first is the initial stress check: a
model that is not in equilibrium before excavation will produce displacements
that have nothing to do with the tunnel, and the check costs one run. The
second is the staging. A single-stage excavation with support installed
simultaneously implies that the lining is in place before any ground movement
occurs, which is never true and which transfers to the lining load that the
ground would have shed. Stress relaxation before support activation, whether
applied as a load-split factor or through a longitudinal deformation profile,
is the two-dimensional representation of the three-dimensional face effect, and
its value is a design decision rather than a default
(Chapter~\ref{ch:ccm}).

\paragraph{Constitutive choice.}
Mohr--Coulomb is the usual first choice and is adequate for soils and for weak
rock at low confinement, where it requires only $c'$ and $\phi'$. Hoek--Brown
is preferable for jointed rock masses at depth, because it captures the
curvature of the strength envelope with confinement that a linear criterion
cannot. The distinction matters less than it appears at shallow depth: over the
confining range of a tunnel at 5 to 25~m, the Hoek--Brown envelope is nearly
straight and equivalent Mohr--Coulomb parameters fitted over that range give
almost the same answer. At several hundred metres the two diverge, and the
linear fit is unconservative in the sidewall and conservative in the crown.

\paragraph{Groundwater.}
Whether to run the analysis dry, with a piezometric surface, or fully coupled
is a modelling decision that follows from the hydrogeology of
Chapter~\ref{ch:hydrogeology} and not from convenience. A dry analysis with
effective stress parameters is appropriate where the water table lies below the
invert or where drainage is rapid. A piezometric surface is appropriate where
the table is known and steady. Neither represents the transient pore pressure
field around an advancing face in low-permeability ground, which requires the
coupled treatment of Section~\ref{sec:coupled-analysis}. Whichever is adopted,
the choice and its basis belong in the design record.

\paragraph{Comparing the model with the empirical estimate.}
The computed trough should be extracted along the ground surface and plotted
against the Gaussian estimate of Section~\ref{sec:gaussian-trough} on common
axes. Table~\ref{tab:comparison} sets out what the comparison can and cannot
establish.

\begin{table}[htbp]
\caption{What a comparison between the empirical and numerical settlement
estimates can establish. Horizontal displacement has no empirical
counterpart.}
\label{tab:comparison}
\centering
\small
\begin{tabularx}{\textwidth}{l c c X}
\toprule
\textbf{Quantity} & \textbf{Empirical} & \textbf{Numerical} & \textbf{Comment} \\
\midrule
Maximum settlement & From $V_L$ and $i$ & Computed & Direct comparison \\
Trough width & $i = K z_0$ & Measured from output & Shape comparison \\
Trough shape & Gaussian assumed & Computed & Overlay the two \\
Horizontal displacement & Not predicted & Computed & Numerical only \\
Governing assumption & $V_L$ and $K$ & Constitutive model and $K_0$ & Different failure modes \\
Principal limitation & Empirical, no stress & Mesh and boundary sensitivity & Both must be stated \\
\bottomrule
\end{tabularx}
\end{table}

The two will not agree exactly and need not be made to.
The empirical curve assumes a volume loss and imposes a shape derived from
soft-ground case histories; the model computes a shape from the stress--strain
response it has been given. Disagreement is informative where its cause can be
named: a layered profile rather than a uniform one, an anisotropic stress
field, a contact near crown level, or cover of less than one diameter. A
difference that cannot be explained is a reason to distrust the model, and
averaging the two does not resolve it.

\subsection{Model extent, mesh and shallow cover}
\label{sec:model-extent-mesh}

Three set-up decisions govern whether a continuum model is numerically sound,
and all three are made before any result exists to judge them by.

\paragraph{Extent.}
Lateral boundaries should stand at least five diameters from the tunnel
centreline, and the base at least five diameters below the invert. Restrained
boundaries placed closer than this stiffen the response and suppress
displacement, which makes the result appear more favourable than it is
(Figure~\ref{fig:extent-mesh}). A shallow tunnel needs a wider domain
than the rule suggests, because the settlement trough is wide relative to the
depth and a boundary within the trough truncates it. The test is empirical:
the boundary is moved outward and the analysis repeated, and a result that
changes materially shows that the boundary was too close.

\paragraph{Mesh.}
A coarse mesh gives inaccurate displacements and structural forces; a uniformly
fine one costs computation without improving either. A graded mesh resolves
this, with elements of roughly 0.3 to 0.5~m within two diameters of the
excavation, one to two metres in the intermediate zone, and three to five
metres in the far field. Mesh adequacy has to be demonstrated: the model is run
at two densities to confirm that the quantity of interest changes by
less than about five per cent. Without that demonstration, mesh dependence
cannot be distinguished from ground behaviour.

\begin{figure}[htbp]
\centering
\datafigure[\textwidth]{ch17/extent-mesh}
\caption{Set-up checked against the closed form for a circular opening unloaded in a hydrostatic field. Left: wall displacement lost to a fixed boundary at radius $R$, 3 per cent at $R = 5D$ for $\nu = 0.3$ but 10 per cent for $\nu = 0.45$; circles are finite element runs. Right: error against element size at the wall, larger for the tangential stress read in the first element than for the displacement.}
\label{fig:extent-mesh}
\end{figure}

\paragraph{Shallow cover.}
The hardest case in this class is a tunnel whose cover is less than its own
diameter. With, say, 1.4~m of ground above the crown of a 7.2~m tunnel, the
plastic zone may reach the surface, displacements become large, and the
solution may fail to converge. The mesh must be fine in a region that is
physically thin, and the answer is sensitive to both the constitutive model and
the stress ratio in a way that it is not at depth.

The numerical remedies are worth trying first: relax the excavation in
increments rather than in a single stage, refine the crown region only rather
than the whole domain, and allow more iterations before declaring
non-convergence. If those fail, further adjustment of the solver is not the useful
response. Non-convergence then indicates that the
ground as idealised cannot stand unsupported. The engineering
response is to model what would be done in construction (ground treatment above the
crown represented as a zone of enhanced strength and stiffness, systematic
radial reinforcement, forepoling or a pipe umbrella represented as a stiffened
crown region, or a stiffer lining) and to record the measure as a design
decision rather than a numerical expedient. Reporting the failure and then
demonstrating the improvement is a stronger result than a converged run
obtained by loosening a tolerance.

\subsection{Scope and limits of numerical analysis}
\label{sec:scope-limits-numerical}

\paragraph{The limits of two-dimensional continuum analysis.}
Plane strain assumes the ground cannot displace along the tunnel axis, which
holds remote from the face and not at it. Since most convergence occurs
within about a diameter of the face, and support is installed there, the
assumption fails where the design decision is made.

Two-dimensional analysis handles this through a stress relaxation factor:
only part of the in-situ stress is released before the support is activated,
the remainder afterwards. The device works, but the answer depends directly on
a factor that the analysis itself cannot determine. The factor is normally
taken from a longitudinal deformation profile, that is, from an axisymmetric or
three-dimensional solution, so the two-dimensional model is calibrated by
the kind of analysis it was intended to replace.

Problem classes where the device fails altogether are those without
longitudinal uniformity: junctions and cross-passages, portals, staged cavern
excavation, and any case where the ground changes over a distance comparable
to the diameter.

\paragraph{Constitutive models beyond Mohr--Coulomb.}
Mohr--Coulomb is the default because it needs two strength parameters and
behaves predictably. Its limitations depend on the problem.

For urban settlement work the deficiency is stiffness, not strength. Soil
stiffness falls sharply with strain, and the small-strain stiffness governing
ground movement outside the yield zone is several times the value measured in
a conventional triaxial test. A Mohr--Coulomb model with a single modulus
gives a settlement trough that is too wide and too shallow, and the trough
shape is the input to a damage assessment. Hardening soil models with
small-strain stiffness address this and are the appropriate choice where
settlement governs.

For brittle rock at high stress the deficiency is the shape of the failure
envelope near the boundary. Cohesion is lost at small strain while friction is
mobilised only at larger strain, so the two are not available simultaneously;
cohesion-weakening friction-strengthening formulations represent this and
predict the shallow spalling notch that Mohr--Coulomb cannot.

For squeezing ground the deficiency is time. Convergence continues for months
after the face has passed, and only creep or viscoplastic formulations
reproduce it.

\paragraph{Selecting and parameterising a model.}
Model complexity should be matched to the data available to feed it. A
hardening soil model requires stiffness parameters at several strain levels; a
creep model requires time-dependent testing; a strain-softening model requires
post-peak behaviour that standard testing does not provide.

Adopting an advanced model and populating it from correlations produces an
analysis whose form is sophisticated but whose content is arbitrary. Where the
parameters have not been measured, a simple model with a sensitivity
study is the better choice.

\section{Methods beyond continuum analysis}
\label{sec:beyond-continuum}

\subsection{The family of numerical methods}
\label{sec:numerical-family}

\paragraph{What each method discretises.}
The methods differ in what they divide into elements, and that difference
determines what each suits (Figure~\ref{fig:method-family}). Finite element and finite difference methods
discretise the whole domain, handling material nonlinearity and heterogeneity
naturally at the cost of meshing a large volume and truncating it at an
artificial boundary. Boundary element methods discretise only surfaces.
Distinct element methods discretise the rock into blocks and solve their
interaction.

\begin{figure}[htbp]
\centering
\datafigure[\textwidth]{ch17/method-family}
\caption{What each family discretises (schematic). The finite element method meshes the whole domain and truncates it at an artificial boundary; the boundary element method discretises only the excavation surface, the far field being exact; the distinct element method divides the rock into blocks, here releasing a key block at the crown.}
\label{fig:method-family}
\end{figure}

\paragraph{Boundary element methods.}
Because only the excavation surface is discretised, the far field is
represented exactly and infinite-medium conditions are satisfied without an
artificial boundary. For elastic problems in homogeneous ground this is both
efficient and accurate, and it remains the quickest way to map the stress
field around a complex three-dimensional excavation geometry.

The restriction is that material nonlinearity and heterogeneity are difficult
to accommodate. Boundary element analysis suits stress distribution studies,
shape comparison and zone-of-influence work in competent ground, and does not
suit yielding, staged construction or layered profiles.

\paragraph{Linked and hybrid schemes.}
Coupling a boundary element far field to a finite or distinct element near
field gives an exact far-field condition with nonlinearity where it occurs.
The approach is less common than its merits suggest, largely
because integrated software support is limited.

\subsection{Discontinuum and hybrid methods}
\label{sec:discontinuum-hybrid}

\paragraph{Discontinuum modelling.}
The distinct element method represents the rock mass as an assembly of blocks
bounded by discontinuities, each block free to translate, rotate and detach,
with contact behaviour governed by joint stiffness and strength. Where block
movement is the failure mechanism, this represents the physics directly rather
than approximating it through an equivalent continuum.

The test is whether explicit joints change the answer or only refine it.
Where the excavation intersects a modest number of blocks and their geometry
governs, they change it. Where the excavation intersects hundreds of blocks
and the mass behaves as a weak continuum, they refine it at considerable cost,
and the equivalent continuum is the appropriate model.

Input requirements are demanding: joint geometry, spacing, persistence and
stiffness in normal and shear. Joint stiffness in particular is rarely
measured and strongly influences the result.

\paragraph{Block theory as a screen.}
Before building a discontinuum model, the removability analysis of
Chapter~\ref{ch:structural-analysis} identifies which blocks can move.
It is quick, deterministic and depends only on orientation, and it frequently
shows that the structural problem is confined to a few block types, which
is either a complete answer or a well-posed brief for the model.

\paragraph{Hybrid and fracture-propagation methods.}
Combined finite--discrete element methods allow intact material to fracture
and the resulting fragments to interact, representing fracture initiation,
propagation and coalescence rather than assuming a pre-existing network. They
are the natural tool for spalling, blast damage and fragmentation, but they
remain research-grade rather than routine: computationally expensive, and
dependent on fracture-energy parameters that few projects measure.

\subsection{Three-dimensional and staged analysis}
\label{sec:three-dimensional-staged}

\paragraph{Three-dimensional modelling.}
Three dimensions are required where geometry or sequence has no longitudinal
uniformity: junctions and cross-passages, portals with sloping cover, cavern
excavation in stages, twin tunnels at close spacing, and fault crossings.
Three-dimensional analysis is also expensive, and only the part of the
problem that needs it should be modelled in three dimensions, with
two-dimensional analysis elsewhere.

Mesh management is the practical constraint. Element size must resolve the
excavation and the yield zone while the domain extends far enough that
boundaries do not influence the result, conventionally three to five
diameters. That relaxes the five-diameter minimum of
Section~\ref{sec:model-extent-mesh}, and the reason is cost: the number of
elements grows roughly with the cube of the domain dimension rather than with
its square. The tighter domain makes the test of moving the boundary outward
and repeating the analysis obligatory. Graded meshes reconcile resolution with extent, and the grading
itself must be checked by refinement.

\paragraph{Simulating construction sequence.}
Staged analysis reproduces the construction: excavate, relax, install support,
advance. Results are sensitive to the sequence assumed, and particularly to
when support is activated and with what stiffness.

Two effects are commonly mishandled. First, shotcrete stiffness rises with curing, so
a lining modelled at mature stiffness from the moment of installation attracts
more load than the real one; young shotcrete also creeps substantially, which
relaxes load. Second, a segmental ring is installed at a defined convergence set
by the tail gap and grouting, so its starting point on the convergence axis is
a construction parameter rather than a modelling choice.

\section{Coupled and multi-physics analysis}
\label{sec:coupled-analysis}

\paragraph{Hydro-mechanical coupling.}
Where pore pressure and deformation interact, they must be solved together.
Excavation changes pore pressure; the changed pressure alters effective stress
and therefore deformation; deformation alters permeability. Consolidation
analysis handles the time-dependent version, steady-state seepage the
long-term one.

Time-dependent settlement from drawdown
(Chapter~\ref{ch:hydrogeology}) requires coupled analysis to predict, as does
long-term load transfer to a lining in low-permeability ground. Drainage
boundary conditions on the lining are the governing input, and
they follow from the drained-or-tanked decision rather than from the analysis.

\paragraph{Thermal, dynamic and other couplings.}
Thermo-hydro-mechanical analysis is required for waste repositories, for
energy geostructures with thermally activated linings, and for tunnels at
depth where rock temperature is significant. Dynamic analysis is required for
blast loading and for seismic response (Chapter~\ref{ch:seismic}), and
introduces input motion selection, boundary conditions that absorb rather than
reflect outgoing waves, and damping, each capable of dominating the result.

\section{Calibration and quality assurance}
\label{sec:calibration-qa}

\paragraph{Calibration, back-analysis and updating.}
Monitoring data allow a model to be calibrated against observed behaviour, and
a calibrated model is a far better predictor for the remainder of a drive than
an uncalibrated one. This is the formal mechanism behind observational design.

The limitation is non-uniqueness. Several combinations of modulus, strength
and in-situ stress reproduce the same measured convergence, and fitting one
combination does not establish that it is correct. Back-analysis is reliable
when it refines parameters already bracketed by measurement, and unreliable
when asked to determine several unknowns from a single observed quantity. The
defence is to calibrate against more than one measured quantity
(convergence, extensometer profile and support load), since a parameter
set reproducing all three is far more constrained.

\paragraph{Verification, validation and quality assurance.}
Verification establishes whether the model solves the equations correctly,
and validation whether the equations describe the problem. Both are
required.

Verification is largely mechanical: mesh convergence studies, comparison
against closed-form solutions for simplified cases, and equilibrium checks.
The closed-form comparison is the most valuable and the most neglected. A
model whose circular-tunnel case does not reproduce the Kirsch solution~\citep{kirsch1898} has an
error that will not be apparent in the complex geometry for which it was built,
and running the check costs an hour (Section~\ref{sec:kirsch-appendix}).

Validation requires evidence external to the model: comparison against
monitoring, against case histories, or against an independent analysis by a
different method. Without any of these, a numerical result remains a hypothesis.

Presenting output as design without stating this distinction is a failure of
professional practice. Results should be reported with the
constitutive model, its parameters and their provenance, the boundary
conditions, the construction sequence assumed, and the verification performed,
so that a reviewer can judge the analysis.

\section{Surrogate models}
\label{sec:surrogate-models}

A numerical model gives one result per run. That is adequate for a design
check but not for tasks that require the result many times over:
reliability analysis, sensitivity study, optimisation, and any twin that must
respond while the excavation is still advancing. A staged three-dimensional
analysis taking hours cannot be placed inside a loop that needs thousands of
evaluations, and the usual response (abandoning the uncertainty analysis
and reporting a single deterministic run) gives a
misleading answer.

The alternative is to run the expensive model a designed number of times, fit
a cheap function to what it produced, and evaluate that function instead. The
cheap function is the surrogate. It does not simplify the physics. It
interpolates the output of one model over one range
of inputs, and it is valid only within those limits.

\subsection{Constructing a surrogate}
\label{sec:constructing-surrogate}

Three choices determine the quality of the surrogate.

The first is which inputs to vary. A surrogate over twenty parameters requires
an impractical number of runs; a surrogate over the three or four that the
sensitivity analysis of Section~\ref{sec:sensitivity-and-project-risk}
identifies as governing is affordable and usually sufficient. The screening
step is therefore a large part of the work.

The second is where to sample. Runs should span the input ranges rather than
cluster near the central estimate, because a reliability calculation evaluates the surrogate
at the tails. Space-filling
designs, of which Latin hypercube sampling is the common
example~\citep{sacks1989design}, distribute a fixed budget of runs across the
input space far more efficiently than varying one parameter at a time, which
samples a cross rather than a volume and cannot detect interactions
(Figure~\ref{fig:surrogate}).

The third is the form of the fit. A quadratic response surface suffices where
the response is smooth and monotonic, which covers a good deal of tunnel
behaviour within the elastic range. A Gaussian process fit is preferable where
the response is curved or where a statement of its own uncertainty is wanted,
since it returns a variance alongside each prediction and that variance grows
away from the sampled points, a property that no polynomial fit provides and that
bears on the misuse described in Section~\ref{sec:surrogate-failure}.

\begin{figure}[htbp]
\centering
\datafigure[\textwidth]{ch17/surrogate}
\caption{Left: nine runs over two inputs, one at a time (squares) and by Latin hypercube (circles). Right: a quadratic and a Gaussian process fitted to five runs of the closed-form convergence of the deep tunnel at GSI 30 over a narrow range of cohesion (shaded). Outside it both fits depart from the parent model; the process band widens, but on the weak side it does not contain the parent.}
\label{fig:surrogate}
\end{figure}

Validation follows the rule that governs all fitted models: hold runs back and
test against them. Reporting the fit quality on the same runs used to fit
demonstrates only that the fitting worked. A surrogate should be reported with
the number of parent runs, the sampled ranges, the held-out error, and the
parameters held fixed.

\subsection{Where surrogates are used}
\label{sec:surrogate-use}

Reliability analysis is the clearest case. Section~\ref{sec:reliability-analysis}
notes that a performance function evaluated by finite element run justifies a
point estimate method or a surrogate; the surrogate is the more informative
choice, because once fitted it supports Monte Carlo sampling, first-order
reliability analysis and sensitivity study at negligible further cost.

Optimisation is the second. Support spacing, lining thickness, relaxation
before installation and excavation sequence trade against one another, and the
trade-offs cannot be explored at an hour per evaluation. The same applies to
the carbon comparisons of Chapter~\ref{ch:sustainability}, where the design
space must be swept rather than sampled at two points.

The third is the digital twin. A twin that predicts runs a model, and one
that predicts fast enough to inform the next round runs a surrogate,
refreshed periodically from the parent model as the ground model is updated.
This is the mechanism that makes the update loop described in
Chapter~\ref{ch:digital-twin} achievable.

\subsection{The failure mode}
\label{sec:surrogate-failure}

A surrogate carries two error terms, and confusing them is the main source of
error. The first is the parent model's error against reality, which the
surrogate carries in full: a surrogate of a model with the wrong constitutive
law reproduces the wrong answer efficiently. The second is the surrogate's own
error against the parent, which is what held-out validation measures. Only the
second is improved by more runs or a better fit, and a report quoting a
held-out error of two per cent is describing agreement with the model, not
with the ground.

The principal hazard is extrapolation. A polynomial surrogate
evaluated outside its sampled range returns a smooth, plausible and unbounded
number with nothing to indicate that it is unsupported
(Figure~\ref{fig:surrogate}), and reliability
calculations reach into the tails where that occurs. Three defences
apply. Sample wider than the range expected to be needed, since the cost is
linear and the failure is silent. Prefer a fit that reports its own predictive
variance, since that variance usually grows away from the sampled points and
so gives some warning of extrapolation, although not a guarantee. Refit rather than
extrapolate when the ground model is updated and the input ranges move.

The final check cannot be automated. A surrogate should be
interrogated at a few points where the answer is known independently (an
analytical solution from Appendix~\ref{app:closed-form}, a hand calculation
from Appendix~\ref{app:calculations}, or a measured convergence), and its
predictions compared. Agreement there is worth more than any goodness-of-fit
statistic, because it tests the parent model and the fit together, where the statistic
tests only the fit.

\chapterreferences
